\begingroup
\fancyhead[L]{\small REMARKS}
\renewcommand{\headrulewidth}{0pt}
\par\noindent\hspace*{\parindent}The following details should be noted.

\par\noindent\hspace*{\parindent}On [I.]: The passages from the letters are printed from the originals,
with the exception of the first letter to \textsc{Olbers} of 3 January 1806,
the original of which is at Pulkovo Observatory. The archive holds a copy
through the mediation of Mr C. \textsc{Schilling}-Bremen.
Concerning the letter from \textsc{Hansen} to \textsc{Gauss} dated 7 February 1843, it should be noted
that it is not \(\dfrac{p''}{p_1''\mu}\) but the mean
\(\frac12\left(\dfrac{p''}{p_1''\mu}+\dfrac{r''}{r_1''\mu}\right)\)
that gives the last-mentioned number \(9{,}791\,4946_n\). Hence it is not
\(\dfrac{p''}{p_1''\mu}\) that agrees with \(\dfrac{p'''}{p_1'''\mu}\)
but this mean. I owe this information to Mr J. \textsc{Kramer},
who recalculated the example.

\par\noindent\hspace*{\parindent}On [II.]: There are two unfinished drafts of the »Exposition d'une nouvelle méthode« etc.
in the posthumous papers. Only the first draft
bears the motto reproduced with the text. The first three »Sections« do not differ
substantially between the two drafts and are printed (pages 439--464) from the second
draft. The »Quatrième Section«, however, is entirely different in the two drafts
and is therefore printed from both. In the first draft it breaks off
in the middle of a sentence. A slip of paper found among the manuscripts bears
the following remark, signed »L[\textsc{isting}] 1855 Aug. 18«:
»This treatise, apparently quite ready for the press, which was found in the \((3+4)\)th volume
of \textsc{Crelle}'s Journal, appears [to] have to be dated to 1806 (at the latest),
inasmuch as the method of expansion in the »Quatrième Section«,
according to \textsc{Gauss}'s oral remarks, was written down earlier than the one
given by \textsc{Fourier} in Dec. 1806«, with the later addition:
»\textls[120]{However}, the passage in Art. 20 »pendant les années 1803 \ldots{} 1816«
places the time of composition after the latter year. It would also be of some interest
to trace the origin of the constants cited at the end of Art. 3.«

\par\noindent\hspace*{\parindent}The correspondence, however, shows that the first
draft in any event must be dated to about 1815, and the second, which deals with the subject mentioned by \textsc{Listing},
therefore to a still later date. The oral remarks by \textsc{Gauss} can refer
only to the fact that he had already carried out these investigations
and may have recorded them elsewhere. From a friendly source I received
some interesting information about \textsc{Gauss}'s proof of convergence.
The only point I can mention here is that \textsc{Poisson}'s analogous proof already
appeared as an extract in 1815 in the Bulletin of the Société philomatique, p. 90.
It is not impossible that \textsc{Gauss} knew of this extract when composing the second
draft and that this very extract may have prompted him to write down
the proof he had found earlier. As already mentioned, \textsc{Gauss} seems
to have undertaken the preparation of the »Exposition« intending to submit it for the
prize offered by the Paris Academy for the perturbations of Pallas.
The award of the prize was ultimately postponed until 1816
(see also the letter to \textsc{Olbers} of 8 January 1816, p. 427).
\endgroup
